
In this section, we first present the main limitations of our method and
outline future directions.

\paragraph{Limitations.} Although \replay{} is more efficient than existing
methods at computing metagradients, it is still non-trivially more expensive
than simply training a model once. The main reason is that metagradients require
making a \emph{backwards pass over a backwards pass}. This operation necessarily
requires 2-3 times the operations of a backwards pass; furthermore, our current
implementation requires \texttt{float32}/\texttt{tensorfloat32} operations.
Finally, standard training operations are often made more efficient by
specialized software (e.g., via FlashAttention~\citep{dao2022flashattention});
no such software (yet) exists for backwards-over-backwards operations. Beyond
computational issues, successfully applying metagradients requires smooth model
training. 

\paragraph{Metasmoothness: connections and future directions.}
While Section \ref{sec:localpred} describes a general procedure for finding 
metasmooth learning algorithms, an important future direction is 
to further explore and understand metasmoothness.
This includes, for example: 
(a) characterizing the relationship between metasmoothness 
and numerical stability (and potentially using techniques from the latter to
improve the former); 
(b) devising improved optimizers and/or architectures that lead directly 
to metasmooth learning algorithms (akin to skip connections or 
stable initialization in architecture design);
(c) formalizing connections between metasmoothness and other 
optimization-related phenomena in deep learning 
\citep{leclerc2020two,cohen2022gradient}. 
A related but separate direction is to explore the possibility of using 
techniques from non-smooth optimization \citep{clarke1990optimization} 
to perform metagradient descent on non-metasmooth learning algorithms.


\paragraph{Applying metagradients.}
Our methods apply to any ML task that requires
optimizing with respect to a metaparameter. 
These include: poisoning data
(generated or simply hosted on the internet) so that it cannot be trained on
without permission (i.e., by maximizing training loss with respect to the text);
selecting better training data at various stages of the model training
lifecycle; and designing better model training routines and architectures with
first-order methods. Another direction of future work lies in mitigating the
computational limitations of our algorithm. Both (a) small-scale
proxy-models~\citep{hoffmann2022training,engstrom2024dsdm} and (b) low-hanging
engineering improvements can likely make calculating metagradients much more
efficient.



